\section{Introduzione}

\subsection{Problem Statement e Obiettivi del Progetto}
La gestione dello smaltimento dei rifiuti urbani ingombranti e speciali nella città di Roma presenta storicamente complessità logistiche legate al coordinamento tra cittadini, centri di raccolta territoriali (isole ecologiche) e squadre operative di raccolta su strada. I canali tradizionali non integrati comportano spesso tempi di attesa elevati per l'utenza, difficoltà nella stima preventiva dei carichi e rischi di saturazione non controllata dei mezzi e delle sedi.

La piattaforma \textbf{MyAma} si propone come soluzione software integrata volta a digitalizzare e ottimizzare l'intero ciclo di vita delle richieste di smaltimento dei rifiuti ingombranti. L'obiettivo primario del sistema è duplice:
\begin{itemize}[leftmargin=*]
    \item \textbf{Per la cittadinanza:} fornire un canale digitale trasparente, accessibile e guidato per prenotare in autonomia sia il \textit{ritiro a domicilio} sia il \textit{conferimento programmato in sede}, consentendo il tracciamento in tempo reale dello stato del servizio ed evitando code e disservizi.
    \item \textbf{Per l'azienda (AMA):} fornire strumenti di pianificazione logistica per l'allocazione ottimizzata dei turni, il rispetto della capacità massima di carico dei veicoli, il contingentamento degli accessi ai centri di raccolta e il monitoraggio puntuale degli esiti operativi.
\end{itemize}

\subsection{Classi di Utenza del Sistema}
Il servizio è accessibile alle seguenti classi di utenza (attori):
\begin{itemize}[leftmargin=*]
    \item \textbf{Utente di sistema:} è qualsiasi utente registrato che ha effettuato con successo l'accesso alla piattaforma. In base al ruolo ricoperto, può accedere esclusivamente alle funzionalità previste dal proprio profilo.
    
    \item \textbf{Utente non registrato:} è un generico individuo che ancora non possiede un account nella piattaforma MYAma. Esso può registrarsi e conseguentemente accedere al profilo personale provvisto di funzionalità specifiche in base al ruolo rivestito nella piattaforma.
    
    \item \textbf{Cittadino:} può consultare liberamente le informazioni generali sui servizi offerti, le tipologie di rifiuti ammesse, le sedi territoriali attive e le relative tariffe. Per procedere alla prenotazione di un servizio, può registrarsi fornendo i propri dati anagrafici e di contatto (nome, cognome, indirizzo, recapito telefonico e indirizzo e-mail). In seguito alla registrazione può usufruire della piattaforma per richiedere un ritiro a domicilio o prenotare un conferimento diretto presso una sede AMA, specificando le caratteristiche del rifiuto (con eventuale caricamento foto), indicando l'indirizzo/CAP e selezionando una fascia oraria disponibile. Può inoltre monitorare lo stato di avanzamento delle proprie richieste, consultare lo storico, rilasciare valutazioni ed eventualmente annullare una prenotazione attiva entro i limiti temporali previsti.
    
    \item \textbf{Autista AMA:} tramite l'applicazione dedicata, consulta l'elenco dei ritiri assegnati per il proprio turno con i dettagli logistici (indirizzo, fascia oraria, tipologia di carico e capienza residua del mezzo), visualizza i recapiti per contattare il cittadino e registra l'esito dell'attività svolta (completato, cittadino assente, rifiuto non conforme).
    
    \item \textbf{Operatore di sede AMA:} gestisce le attività di accettazione presso il centro di raccolta, verificando le prenotazioni dei cittadini in arrivo, controllando la conformità dei rifiuti conferiti e registrando l'esito del servizio.
    
    \item \textbf{Amministratore di sede AMA:} gestisce l'organizzazione logistica della propria struttura: genera i codici di invito per il personale operativo (autisti e operatori) della sede, definisce le disponibilità di lavoratori e veicoli, imposta le fasce orarie e associa le sedi alle rispettive zone o CAP serviti.
    
    \item \textbf{Amministratore generale AMA:} opera a livello direttivo aziendale; è responsabile della gestione degli account degli Amministratori di sede (generazione codici di invito dedicati, abilitazione e revoca).
\end{itemize}

\subsection{Perimetro del Sistema (In-Scope e Out-of-Scope)}
Al fine di delimitare con precisione i requisiti e i modelli OOA sviluppati nel presente documento:
\begin{itemize}[leftmargin=*]
    \item \textbf{Funzionalità In-Scope:} autenticazione e registrazione basata su ruoli; gestione completa delle prenotazioni di ritiro a domicilio e conferimento in sede; verifica territoriale vincolata a sede e CAP; gestione delle disponibilità e della capacità dei mezzi; tracciamento e registrazione degli esiti di servizio; consultazione dello storico e rilascio feedback; gestione amministrativa di sedi, personale, veicoli e codici di invito.
    \item \textbf{Funzionalità Out-of-Scope (future estensioni):} integrazione diretta con gateway di pagamento elettronico per tariffe extra-franchigia (gestite in questa fase a livello di preventivo informativo); autenticazione federata SPID/CIE; tracciamento GPS in tempo reale dei veicoli; gestione di reportistica analitica avanzata di Business Intelligence.
\end{itemize}
